\documentclass[11pt,a4paper]{article}
\usepackage[a4paper,margin=2.4cm]{geometry}
\usepackage{booktabs}
\usepackage{longtable}
\usepackage{array}
\usepackage{amsmath}
\usepackage{float}
\usepackage{hyperref}
\hypersetup{hidelinks}
\setlength{\emergencystretch}{3em}

\title{\textbf{Supplementary Material:\\
Time-of-Use Pricing for Fixed-Capacity Inference Services}}
\author{Anonymous Author}
\date{}

\begin{document}
\maketitle

\section*{S1. Parameters}

Table~\ref{tab:params} reports the main parameters used in the peak-shaving
experiments. Unless stated otherwise, values come from
\path{pricing_sim/peak_shaving_config.py}. These are synthetic calibrations for
mechanism analysis, not production estimates.

\begin{longtable}{p{0.28\textwidth}p{0.23\textwidth}p{0.39\textwidth}}
\caption{Main simulation parameters}
\label{tab:params}\\
\toprule
Parameter & Value & Meaning \\
\midrule
\endfirsthead
\toprule
Parameter & Value & Meaning \\
\midrule
\endhead
$T$ & 8 & Number of periods in one day. \\
$h$ & 3 hours & Length of one period. \\
$G_A,G_B$ baseline & 1500, 600 & Fixed provider capacities. \\
$G_A,G_B$ congested & 300, 120 & Tightened capacities used to activate QoS degradation. \\
$\pi_R,\pi_E$ baseline & 0.6, 0.4 & Population shares of time-rigid and time-flexible users. \\
$\pi_R,\pi_E$ congested & 0.4, 0.6 & Population shares in the congested experiments. \\
\texttt{base\_rigid} & \begin{tabular}[t]{@{}l@{}}200,150,300,600,\\800,900,700,250\end{tabular} & Rigid-demand baseline. \\
\texttt{time\_preference} & \begin{tabular}[t]{@{}l@{}}0.30,0.20,0.50,0.80,\\1.00,1.00,0.90,0.40\end{tabular} & Exogenous period preference. \\
\texttt{intermediary\_brand} & 1.05 & Intermediary convenience parameter. \\
\texttt{firm\_brand} & 1.0, 1.0 & Providers are nearly homogeneous in model quality. \\
\texttt{flexible\_baseline} & 400 & Flexible-demand scale. \\
\texttt{rigid\_wtp} & 1.80 & Reservation value for rigid or replay demand. \\
\texttt{rigid\_churn\_rate} & 5.0 & Price-response slope for rigid demand. \\
\texttt{market\_growth} & 1.2 & Flexible market expansion at low prices. \\
$p_0$ & 0.80 & Baseline price. \\
Direct price bounds & [0.45, 2.10] & Provider direct retail-price bounds. \\
Wholesale price bounds & [0.25, 0.90] & Provider wholesale-price bounds. \\
$\alpha_R,\alpha_E$ & 2.0, 5.0 & Price sensitivity of time-rigid and time-flexible users. \\
$c_s^R,c_s^E$ & 0.60, 0.10 & Switching costs of time-rigid and time-flexible users. \\
$\bar u$ & 0.82 & QoS degradation threshold. \\
\texttt{qos\_strength} & 15.0 & QoS degradation strength. \\
$\gamma$ & 1.0 & QoS feedback weight in utility. \\
$c_g$ & 0.015 & Fixed capacity holding cost. \\
$c_q$ & 0.35 & QoS degradation cost. \\
$u_0$ & 0.0 & Utility of the no-purchase outside option. \\
\bottomrule
\end{longtable}

\section*{S2. Pricing Shape and Grids}

Provider strategies are restricted to four scalar parameters,
$(\bar w_m,\delta_m,\bar p^D_m,\delta^D_m)$. The period prices are
\[
  w_{m,t}=\bar w_m(1+\delta_m\widehat{\ell}_t),\qquad
  p^D_{m,t}=\bar p^D_m(1+\delta^D_m\widehat{\ell}_t),
\]
with clipping to the allowed price bounds. Uniform pricing is implemented by
setting $\widehat{\ell}_t=0$, so internal $\delta$ values do not create
period-varying prices.

The coarse fictitious-play grid uses
\[
\begin{aligned}
\bar w&\in\{0.30,0.85\},&
\delta&\in\{-0.4,-0.133,0.133,0.4\},\\
\bar p^D&\in\{0.70,1.50\},&
\delta^D&\in\{-0.4,-0.133,0.133,0.4\}.
\end{aligned}
\]
The fine-grid check uses
\[
\begin{aligned}
\bar w&\in\{0.27,0.575,0.88\},&
\delta&\in\{-0.4,-0.2,0,0.2,0.4\},\\
\bar p^D&\in\{0.60,1.10,1.60\},&
\delta^D&\in\{-0.4,-0.2,0,0.2,0.4\}.
\end{aligned}
\]
The intermediary response is computed over a three-parameter grid: retail base
price, retail dynamic intensity, and routing price sensitivity.

\section*{S3. Artifact Map}

\begin{longtable}{p{0.28\textwidth}p{0.31\textwidth}p{0.31\textwidth}}
\caption{Artifact sources for main-paper quantities}
\label{tab:artifact-map}\\
\toprule
Main-paper quantity & Artifact & Fields \\
\midrule
\endfirsthead
\toprule
Main-paper quantity & Artifact & Fields \\
\midrule
\endhead
Uncongested profit and utilization & \path{artifacts/peak_shaving/20260618/peak_shaving_summary.json} & \path{uniform}, \path{dynamic}. \\
Uncongested demand shift & \path{artifacts/peak_shaving/20260618/peak_shaving_summary.json} & \path{P1_migration.rigid_shift}, \path{P1_migration.elastic_shift}. \\
Elastic-population scan & \path{artifacts/peak_shaving/20260618/peak_shaving_summary.json} & \path{elastic_sweep}. \\
Congested uniform baseline & \path{artifacts/peak_shaving/20260618/peak_shaving_congested_fp.json} & \path{uniform.system_profit}, \path{uniform.peak_util}, \path{uniform.min_qos_firm}. \\
Dynamic coarse grid & \path{artifacts/peak_shaving/20260618/peak_shaving_fp_dynamic_converged.json} & \path{system_profit}, \path{peak_util}, \path{min_qos}, \path{delta_A}, \path{delta_B}. \\
Dynamic fine snapshot & \path{artifacts/peak_shaving/20260618/peak_shaving_fp_dynamic_converged_fine.json} & \path{system_profit}, \path{peak_util}, \path{min_qos}, \path{delta_A}, \path{delta_B}. \\
Fine-grid pure-snapshot regret & \path{artifacts/peak_shaving/20260618/fp_dynamic_ckpt_fine.json} & \path{round=40}, \path{maxreg=22.3029}. \\
Coarse-grid regret & \path{artifacts/peak_shaving/20260618/fp_dynamic_ckpt.json} & \path{round=22}, \path{maxreg=4.9821}. \\
Mechanism diagnostics & \path{artifacts/peak_shaving/20260619/} & \path{peak_shaving_diagnostics.json}, CSV summaries, time profiles. \\
Mixed-strategy diagnostic & \path{artifacts/peak_shaving/20260619_submission/peak_shaving_mixed_oracle.json} & \path{trace}, \path{expected_metrics}, \path{row_mix}, \path{col_mix}. \\
Parameter stress tests & \path{artifacts/peak_shaving/20260619_submission/peak_shaving_parameter_sweep_summary.json} & QoS gains, peak-utilization reductions, profit signs. \\
SMPT diagnostics & \path{artifacts/peak_shaving/20260619_smpt/} & Strong baselines, ablations, phase grid, residuals, restricted re-solve. \\
Public price scale & \path{artifacts/peak_shaving/20260619_submission/public_api_price_anchor.json} & Z.AI, DeepSeek, and OpenAI prices checked on 2026-06-19. \\
vLLM QoS anchor & \path{artifacts/peak_shaving/20260619_submission/vllm_qos_anchor_summary.json}, \path{vllm_qos_anchor_points.csv} & TTFT-SLA curves, fitted parameters, and source CSVs. \\
\bottomrule
\end{longtable}

\section*{S4. Submission-Strengthening Diagnostics}

The high-regret fine-grid pure snapshot could be misread as an equilibrium if
reported alone. The script \path{experiments/run_peak_shaving_mixed_oracle.py}
uses a double-oracle diagnostic over the full 225-point fine grid. It
starts from a small support, repeatedly adds best deviations, and solves an
empirical fictitious-play profile on the restricted support. The current run
obtains a $5\times5$ support in round 5. Full-grid maximum regret is $0.2025$.
The expected peak utilization is $0.7030$, the minimum QoS is $0.9699$, and
system profit is $1733.13$.

The script \path{experiments/run_peak_shaving_parameter_sweep.py} performs the
local fixed-policy stress test. It evaluates nine perturbation scenarios covering
capacity scale, price sensitivity, switching cost, and the QoS threshold. Both
dynamic snapshots keep positive minimum-QoS gains and positive peak-utilization
reductions in all nine scenarios. Profit is not stable: the dynamic fine snapshot
is positive in only one scenario.

\section*{S5. SMPT Additional Diagnostics}

The script \path{experiments/run_peak_shaving_smpt_experiments.py} adds diagnostics
requested by the SMPT-oriented review. Outputs are stored under
\path{artifacts/peak_shaving/20260619_smpt/} and
\path{figures/peak_shaving_smpt/}. Table~\ref{tab:smpt-files} lists the generated
files.

\begin{longtable}{p{0.32\textwidth}p{0.55\textwidth}}
\caption{SMPT diagnostic artifact bundle}
\label{tab:smpt-files}\\
\toprule
File & Contents \\
\midrule
\endfirsthead
\toprule
File & Contents \\
\midrule
\endhead
\path{peak_shaving_smpt_experiments.json} & Full bundle with metadata, baselines, ablations, phase grid, fixed-point residuals, and restricted re-solved sensitivity. \\
\path{smpt_baselines.csv} & Static pricing, off-peak discount only, peak surcharge only, dynamic coarse/fine, equal routing, and admission-control diagnostic rows. \\
\path{smpt_ablations.csv} & Outside-option, direct-channel, capacity-homogeneity, time-flexibility, and equal-routing ablations. \\
\path{smpt_phase_grid.csv} & Fixed-policy grid over capacity scale and price-sensitivity scale. \\
\path{smpt_fixed_point_residuals.csv} & Iterations and final residuals for the main baseline and dynamic-policy fixed points. \\
\path{smpt_vv_damping_initial.csv} & Initial-condition and damping-factor audit for uniform, dynamic coarse, and dynamic fine fixed points. \\
\path{smpt_resolved_sensitivity.csv} & Restricted local candidate re-solve for five selected perturbation scenarios. \\
\bottomrule
\end{longtable}

The strongest new baseline comparison is not a new equilibrium proof. It shows
that one-sided price shapes improve QoS but do not match the full dynamic coarse
policy: off-peak-only discounting reaches minimum QoS $0.833$, peak-only
surcharge reaches $0.921$, and dynamic coarse reaches $0.968$. Equal routing
with the same dynamic coarse prices reaches $0.881$, showing that routing and
price timing both contribute. The admission-control diagnostic indicates that
matching the dynamic coarse peak-utilization target with static pricing would
require a minimum admitted fraction of $0.902$ in the tightest period.

The fixed-point residual check reports final residuals below $10^{-9}$ for the
static baseline, off-peak-only, peak-only, dynamic coarse, and dynamic fine rows.
The damping and initialization audit evaluates three initial QoS levels and
three damping factors for uniform, dynamic coarse, and dynamic fine. It converges
in 24 of 27 rows. Dynamic coarse and dynamic fine converge in all 18 rows; the
three non-converged rows are the static uniform baseline with damping $0.5$.
The fixed-policy phase grid has positive QoS gain and peak-utilization reduction
in all 25 capacity--elasticity grid points. The restricted local candidate
re-solve keeps positive QoS gain in all five selected scenarios; profit is
positive in three and negative in two, preserving the profit-boundary conclusion.

\section*{S6. Price and QoS Anchors}

The public price anchor is stored in
\path{artifacts/peak_shaving/20260619_submission/public_api_price_anchor.json}.
It is used only to check order of magnitude. It is not an estimate of demand
elasticity.

The controlled vLLM QoS-shape anchor is produced by
\path{experiments/build_peak_shaving_measurement_anchor.py}. Inputs are
\path{artifacts/vllm-study/20260531-190126/controlled_aggregate.csv} and
\path{artifacts/vllm-study-qwen25-3b/20260531-214710/controlled_aggregate.csv}.
Both scans use vLLM 0.22.0 and an NVIDIA GeForce RTX 4090. QoS is the
attainment rate for a 0.5-second time-to-first-token SLA. Table~\ref{tab:vllm}
reports the fitted threshold-type QoS parameters.

\begin{table}[H]
\centering
\caption{Controlled vLLM single-GPU QoS-shape anchor}
\label{tab:vllm}
\footnotesize
\begin{tabular}{@{}lrrrrrr@{}}
\toprule
Profile & Healthy conc. & $\bar u$ & $\kappa$ & RMSE & Drop conc. & Min QoS \\
\midrule
Qwen2.5-0.5B & 224 & 1.000 & 0.517 & 0.068 & 384 & 0.500 \\
Qwen2.5-3B & 224 & 1.058 & 0.942 & $6.3\times10^{-10}$ & 384 & 0.667 \\
\bottomrule
\end{tabular}
\end{table}

\section*{S7. Reproduction Commands}

Run the following commands in
\path{/root/paper_code/0427_tokenrl/paper_token_cross_survey}. The public
repository is \url{https://github.com/cccht/paper_token_price}; a formal
submission should cite a frozen archival release when available.

\begin{verbatim}
uv run --with-requirements requirements.txt \
  python experiments/build_peak_shaving_diagnostics.py
uv run --with-requirements requirements.txt \
  python experiments/run_peak_shaving_parameter_sweep.py
uv run --with-requirements requirements.txt \
  python experiments/run_peak_shaving_mixed_oracle.py
uv run --with-requirements requirements.txt \
  python experiments/run_peak_shaving_smpt_experiments.py
uv run --with-requirements requirements.txt \
  python experiments/run_peak_shaving_smpt_vv_audit.py
uv run --with-requirements requirements.txt \
  python experiments/build_peak_shaving_measurement_anchor.py
\end{verbatim}

Compile the English manuscript and supplement with:

\begin{verbatim}
xelatex -interaction=nonstopmode -halt-on-error \
  peak_shaving_dynamic_pricing_sci_en_2026-06-19.tex
bibtex peak_shaving_dynamic_pricing_sci_en_2026-06-19
xelatex -interaction=nonstopmode -halt-on-error \
  peak_shaving_dynamic_pricing_sci_en_2026-06-19.tex
xelatex -interaction=nonstopmode -halt-on-error \
  peak_shaving_dynamic_pricing_sci_en_2026-06-19.tex
xelatex -interaction=nonstopmode -halt-on-error \
  peak_shaving_dynamic_pricing_sci_en_supplement_2026-06-19.tex
\end{verbatim}

\section*{S8. Known Boundaries}

The fine-grid pure snapshot does not reach the regret $<5$ criterion within
40 rounds. Its profit, dynamic intensity, and demand-centroid values should be
read as high-regret snapshot diagnostics. The mixed-strategy diagnostic lowers
regret on the finite 225-point grid, but it is not a continuous-strategy
equilibrium proof.

The SMPT diagnostics strengthen the baseline and sensitivity evidence but do not
remove the calibration boundary. The phase grid is still fixed-policy evidence,
and the re-solved sensitivity uses a restricted local candidate set rather than
a full continuous-strategy search. The paper does not calibrate price sensitivity,
switching cost, or outside-option utility from real request traces. Public API
prices provide only a price-scale anchor. The vLLM scans provide only a single-GPU,
small-model, short-generation QoS-shape anchor. Production use would require
traces, multi-GPU measurements, long-context workloads, user behavior calibration,
and online evaluation.

\end{document}
